\documentclass[10pt,a4paper]{amsart}
\usepackage[margin=19mm]{geometry}
\usepackage{amsmath,amssymb,mathtools}
\usepackage[scr=rsfs,cal=euler]{mathalfa}
\usepackage{xfrac}
\usepackage[unicode,hidelinks,bookmarks=false]{hyperref}
\newcommand{\ie}{i.e.\ }
\newcommand{\cf}{cf.\ }
\newcommand{\resp}{respectively\ }
\newcommand{\Subsection}{\subsection}
\providecommand{\nobreakdash}{\nobreak\hbox{-}}
\setlength{\emergencystretch}{2em}
\multlinegap0pt
\usepackage{bm}
\usepackage{bbm}
\usepackage{dsfont}
\usepackage{xspace}
\usepackage{cancel}
\usepackage{mathtools}
\usepackage[shortlabels]{enumitem}
\usepackage{amsthm}
\makeatletter
\@ifundefined{lemma}{\newtheorem{lemma}{Lemma}}{}
\@ifundefined{theorem}{\newtheorem{theorem}{Theorem}}{}
\@ifundefined{proposition}{\newtheorem{proposition}[theorem]{Proposition}}{}
\makeatother
\newcommand{\R}{\mathbb R}
\newcommand{\mc}{\mathcal}
\title[A Minimal Dynamical Model for Incubation-Outbreak Transition]{A Minimal Dynamical Model for Incubation-Outbreak Transitions in Social Norm Diffusion}
\author{Run Wang, Leonardo Cianfanelli, Giacomo Como, Wenjun Mei}
\date{Source: arXiv:2607.25586v1 (CC BY 4.0)}
\begin{document}
\maketitle

\noindent\emph{Source and licence.} A Minimal Dynamical Model for Incubation-Outbreak Transitions in Social Norm Diffusion. Version arXiv:2607.25586v1 (CC BY 4.0), distributed under the Creative Commons Attribution 4.0 International licence, as recorded in the arXiv licence metadata for this exact version and shown on the abstract page https://arxiv.org/abs/2607.25586v1. Full rights notice: licenses/arxiv-2607.25586-CC-BY-4.0.txt.\par\smallskip

\noindent\textit{Editorial scope.} Counted unit: the Theorem whose source label is \texttt{None} in the article (source0.tex lines 380--386, source label \texttt{None}), delivered together with the article's own complete original proof of it (source0.tex lines 388--436, one \texttt{proof} environment of 49 lines). No mathematics is written, repaired or re-derived: statement and proof are the article's own text line for line. Required context retained uncounted: sys\_dynamics (lines 92--98); proposition:basic (lines 112--118); prop:equilibria (lines 130--149); prop:H (lines 238--240); lemma:g (lines 283--317). Every definition, display and label of those lines is retained unchanged. Omitted from this package and not redistributed: 1--91, 99--111, 119--129, 150--237, 241--282, 318--379, 387--387, 437--684. Nothing inside a retained excerpt is omitted or altered: every retained body is delivered exactly as the article's lines are written, so no omission receipt of any kind accompanies this package. Citations: the retained excerpts carry no citation command at all, so no bibliography entry of the article is transcribed and no reference is redistributed. Reference adaptation: none was needed and none was made; every reference inside the delivered unit resolves inside it, so no unresolved reference or citation is produced. Printed slips kept literal: no printed word, symbol, hypothesis or number of the counted statement or of its proof was altered. Layout and class: the article's own class is \texttt{article} with packages including babel, geometry, amsmath, graphicx, float, tikz, amssymb, caption, subcaption, hyperref, amsthm, mathtools, enumitem, indentfirst; the wrapper is the lane's pinned \texttt{amsart} editorial wrapper at 10pt on A4, which restates the theorem-like environments the retained text needs and adds the article's own macro definitions that the retained text needs (\texttt{\textbackslash R}, \texttt{\textbackslash mc}), with extra wrapper packages (bm, bbm, dsfont, xspace, cancel, mathtools, enumitem). The delivered numbering is the unit's own. Assets: no retained span contains a figure environment, a graphics-inclusion command, a PDF-page inclusion, a table or a TikZ picture; no image file is copied into this package, the package declares no asset dependency and no image file is redistributed with it. Licence: version arXiv:2607.25586v1 of this article is distributed under the Creative Commons Attribution 4.0 International licence, as recorded in the arXiv licence metadata for this exact version and shown on the abstract page https://arxiv.org/abs/2607.25586v1. A full rights notice is delivered at licenses/arxiv-2607.25586-CC-BY-4.0.txt. Characters: every retained excerpt is pure ASCII, so the delivered engine, XeLaTeX with its default Latin Modern text font, reports no missing character.
\begin{equation}
\label{sys_dynamics}
\begin{cases}
    \dot{x} = (1 - x)y\,, \\
    \dot{y} = \beta(x - y)y - \gamma y\,,
\end{cases}
\end{equation}

\begin{proposition}\label{proposition:basic}
Consider model \eqref{sys_dynamics}. Then: 
\begin{enumerate}
    \item[(i)] for every initial condition $(x(0), y(0))$ in $A$, there exists a unique solution $(x(t), y(t))$. The solution is $\mc C^\infty$ and belongs to $A$ for every $t$ in $\R_+$;
    \item[(ii)] $x(t)$ is monotonically non-decreasing.
\end{enumerate}
\end{proposition}

\begin{proposition}\label{prop:equilibria}
\begin{enumerate}
\item[(i)] If $\beta \le \gamma$, the set of equilibrium points of \eqref{sys_dynamics} in $A$ is given by
\begin{equation}\label{eq:stable_below}
\{(x,0) \in A: x \in [0,1]\}\,,
\end{equation}
and all equilibria are stable but not asymptotically stable. 
\item[(ii)] If $\beta > \gamma$, the set of equilibrium points of \eqref{sys_dynamics} in $A$ is given by 
$$\{(x,0) \in A: x \in [0,1]\} \cup \{(1,y^*)\}\,, \quad y^* = 1-\frac{\gamma}{\beta}\,.$$
Moreover, the equilibrium points 
\begin{equation}\label{eq:stable}\Big\{(x,0) \in A: x \in \Big[0,\frac\gamma\beta\Big)\Big\}\end{equation}
are stable but not asymptotically stable, the equilibrium points
\begin{equation}\label{eq:unstable}\Big\{(x,0) \in A: x \in \Big[\frac\gamma\beta,1\Big]\Big\}\end{equation}
are unstable, and the equilibrium point
\begin{equation}\label{eq:as_stable}
(1,y^*)
\end{equation}
is asymptotically stable.
\end{enumerate}
\end{proposition}

\begin{proposition}\label{prop:H}
    $H_{\beta}(x,y)$ is an invariant of motion of model \eqref{sys_dynamics}, i.e., $\dot H_\beta(x,y) = 0$.
\end{proposition}

\begin{lemma}\label{lemma:g}
Let $\beta > \gamma$. Then: 
\begin{enumerate}
    \item[(i)] given $\varepsilon \ge 0$, for every $(x,y)$ in $A$ with $x < 1$, $H_\beta(x,y) = H_\beta(\gamma/\beta,0) + \varepsilon$ if and only if $y = g_\varepsilon(x)$;

    \item[(ii)] if $\beta \ge 1$, $g_\varepsilon(x)$ is strictly convex for every $\varepsilon \ge 0$. If $\beta < 1$, let
    $$
    \varepsilon_\beta = \frac{(1-\gamma/\beta)^{1-\beta}}{1-\beta} = -H_\beta(\gamma/\beta,0)>0\,.
    $$
Then, $g_\varepsilon(x)$ is strictly convex if $\varepsilon< \varepsilon_\beta$, affine if $\varepsilon = \varepsilon_\beta$, and strictly concave if $\varepsilon > \varepsilon_\beta$;

    \item[(iii)]
    for every $\varepsilon \ge 0$ such that $g_\varepsilon(x)$ is convex, $g_\varepsilon(x)$ always admits a unique stationary point $(\hat x_\varepsilon,\hat y_\varepsilon)$ in $A$ given by
    \begin{equation}\label{eq:xy}    
        \hat{x}_{\varepsilon} =\begin{cases}
        1 - (1 - \gamma)e^{- \varepsilon}, & \beta = 1\\
        %1 - \displaystyle\frac{1}{\big((1 - \gamma/\beta)^{1 - \beta}-\varepsilon(1 - \beta)\big)^{1/(\beta - 1)}} 
        1-((1-\beta)C_{\varepsilon,\beta})^{1/(1-\beta)}
        & \beta \neq 1\,,
        \end{cases}\qquad \hat{y}_{\varepsilon} = \hat x_\varepsilon - \frac{\gamma}{\beta}\,.
\end{equation}
Moreover, for every $\varepsilon \ge 0$
\begin{equation}\label{eq:mono}
\frac{\partial \hat x_\varepsilon}{\partial \varepsilon} = \frac{\partial \hat y_\varepsilon}{\partial \varepsilon} > 0\,;
\end{equation}
        \item[(iv)] for every $\varepsilon \ge 0$ such that $g_{\epsilon}$ is convex,
        \begin{equation*}
            \hat y_\varepsilon \leq \Big(1 - \displaystyle\frac{\gamma}{\beta}\Big)^{\beta}\varepsilon\,;
        \end{equation*}
        \item[(v)] for a sufficiently small $\varepsilon \ge 0$, there exists finite $K_\varepsilon>0$ such that
        $$
        g_{\varepsilon}(x) \leq g_{\varepsilon}(\hat{x}_{\varepsilon}) + K_\varepsilon (x-\hat{x}_{\varepsilon})^2\,, \quad \forall x \in [\hat{x}_{\varepsilon} - \sqrt{\varepsilon}, \hat{x}_{\varepsilon} + \sqrt{\varepsilon}]\,.
        $$
    \end{enumerate}
\end{lemma}

\begin{theorem}
    Let $\beta > \gamma$. Then, for every initial condition $(x(0),y(0)) = (x_0,y_0)$ in $A$:
    \begin{enumerate}[label=(\roman*)]
        \item if $(x_0, y_0)\in A_s$, then $$(x(t), y(t)) \xrightarrow{t \to + \infty} (x^*, 0)\,,$$ where $x^*$ is the smallest solution of $H_\beta(x^*, 0) = H_\beta(x_0, y_0)$ if $y_0>0$, and $x^* = x_0$ if $y_0 = 0$.
        \item If $(x_0, y_0) \in A \setminus A_s$, then $$(x(t), y(t)) \xrightarrow{t \to + \infty} (1, y^*)\,, \qquad y^* = \displaystyle 1-\frac{\gamma}{\beta}\,.$$
    \end{enumerate}
\end{theorem}

\begin{proof}
Since $x(t)$ is monotonically non-decreasing by Proposition \ref{proposition:basic}(ii), then $x(t)$ admits a limit $x^*$ for every initial condition $(x_0,y_0)$. 

(i) If $y_0 = 0$, then $(x_0,y_0)$ is an equilibrium and the claim is trivial. Assume now that $y_0>0$. If $(x_0,y_0)$ belongs to $A_s$, we have that $y_0 < g_0(x_0)$, and $(x(t),y(t))$ can leave $A_s$ only by crossing the curve $y = g_0(x)$, which is prevented by Proposition \ref{prop:H} and Lemma \ref{lemma:g}(i) with $\varepsilon = 0$. If instead $(x_0,y_0)$ lies on the curve $y = g_0(x)$, then $(x(t),y(t))$ moves along the curve until convergence to the equilibrium point $(\gamma/\beta,0)$. This yields $x^* \le \gamma / \beta$. It then follows from Proposition \ref{proposition:basic}(ii) that $x(t) \le \gamma/\beta$ for every $t \ge 0$, hence, from \eqref{sys_dynamics}, $\dot y(t) \le 0$. Therefore, $(x(t),y(t))$ must converge to an equilibrium point $(x^*,0)$ in $A_s$. In particular, since $x(t)$ is continuous and non-decreasing in time, and since every element $(x,0)$ is an equilibrium point, $x^*$ is the smallest element greater than $x_0$ such that $H_\beta(x^*,0) = H_\beta(x_0,y_0)$ due to Proposition \ref{prop:H}.

(ii) Consider the sets $$C = \{(x,y) \in A: y > g_0(x)\}, \quad D = \{(x,y) \in A: 0 < y \le g_0(x), \ x > \gamma/\beta\}$$
and note that $A \setminus A_s = C \cup D$. We now prove that every initial condition in $C$ and $D$ converges to $(1,x^*)$. 

Consider $(x_0, y_0)$ in $C$. Since the trajectory $(x(t),y(t))$ cannot intercept the curve $y = g_0(x)$, then $C$ is positively invariant. We now argue by contradiction. Assume that $x^* \le \gamma /\beta$. By the same arguments used to prove item (i), this would imply that the limit point is in the form $(x^*,0)$, contradicting the fact that $C$ is positively invariant. Hence, $x^* > \gamma /\beta$.
Fix now a small $\delta>0$ such that
$$
x^*-\delta>\frac{\gamma}{\beta}.
$$
Since $x(t)\to x^*$, there exists $T_\delta>0$ such that
$$
x^* - \delta \le x(t) \le x^*\,,
\qquad \forall t\ge T_\delta\,.
$$
Hence, for every $t\ge T_\delta$,
$$
\big(\beta(x^*-\delta-y(t))-\gamma\big)y(t) \le \dot y(t)
\le
\big(\beta(x^*-y(t))-\gamma\big)y(t).
$$
It thus follows that
$$
x^*-\delta-\frac{\gamma}{\beta}
\le
\liminf_{t\to+\infty}y(t)
\le
\limsup_{t\to+\infty}y(t)
\le
x^*-\frac{\gamma}{\beta}.
$$
Letting $\delta\to0$, we conclude that
$$
\lim_{t \to + \infty} (x(t),y(t)) = \Big(x^*,x^*-\frac{\gamma}{\beta}\Big).$$
Since the only equilibrium point with positive second component is $(1,1-\gamma/\beta),
$
we conclude that $x^*=1$, and therefore
$$
\lim_{t \to +\infty} (x(t),y(t))
=
\left(1,1-\frac{\gamma}{\beta}\right) = (1,y^*)\,,
$$
concluding the proof of the case $(x_0,y_0)$ in $C$.

Finally, consider the case where $(x_0,y_0)$ belongs to $D$. Note that $D$ is positively invariant, and $\dot x,\dot y \ge 0$ for every $(x,y)$ in $D$. This proves that $x(t),y(t)$ admits a limit point and such limit point must be $(1,y^*)$ by Proposition \ref{prop:equilibria}.
\end{proof}

\end{document}
